\documentclass[12pt]{article}
\usepackage{amsmath,amssymb,mathtools}
\usepackage{hyperref}

\DeclareMathOperator{\Res}{Res}
\DeclareMathOperator{\Ad}{Ad}
\DeclareMathOperator{\vol}{vol}
\newcommand{\A}{\mathbb{A}}
\newcommand{\calW}{\mathcal{W}}
\newcommand{\calO}{\mathcal{O}}

\title{Global Residue Positivity at $s=1$\\
\large Proof Document 02 --- sym2-effective-lbound}
\author{}
\date{}

\begin{document}
\maketitle

\begin{abstract}
This document records the open obligation F-2.  The desired statement is that a
suitably normalized global Jacquet--Shalika quadratic integral has a strictly
positive residue at $s=1$.  The earlier version incorrectly promoted this
statement to a theorem while two inputs were not source-backed: an exact
$L(1,\pi,\Ad)>0$ theorem was not located, and the positivity and normalization
of every ramified/archimedean local correction factor were not proved.  F-2 is
therefore \texttt{[OBL]} and may not be used downstream.
\end{abstract}

\section*{Status}

\textbf{Obligation F-2: \texttt{[OBL].}}  The statement below is not a theorem
in this repository.  In particular, it does not imply any lower bound for
$L(1,\mathrm{sym}^2 f)$, any quantitative estimate uniform in the level, or the
absence of a Siegel zero.

\section{Desired statement}

Let $\pi=\bigotimes_v'\pi_v$ be an irreducible generic cuspidal automorphic
representation of $GL_2(\A_F)$ with contragredient $\tilde\pi$.  Fix a nonzero
additive character $\psi:F\backslash\A_F\to\mathbb{C}^\times$, a factorizable
Schwartz--Bruhat function $\Phi$, and nonzero factorizable Whittaker vectors
$W\in\calW(\pi,\psi)$ and $\tilde W\in\calW(\tilde\pi,\psi^{-1})$.  The desired
global object is a Jacquet--Shalika integral of the form
\[
  \Psi(s,W,\tilde W,\Phi)
  =\int_{N_2(\A_F)\backslash GL_2(\A_F)}
    \tilde W(g)W(g)\Phi(e_2g)|\det g|_{\A_F}^{s}\,d^\times g .
\]
The obligation is to prove, with all normalizations explicit, that after the
standard Rankin--Selberg regularization it has a simple pole at $s=1$ whose
residue is strictly positive for the intended class of real data, and that this
residue factors as a strictly positive constant times $L(1,\pi,\Ad)$.

\section{Source-backed inputs and remaining gaps}

\subsection{Unramified local integral}

For generic unramified local data, Jacquet--Shalika define the essential
Whittaker functions and prove
\[
  \Psi_v(s,W^\circ,\widetilde W^\circ,\mathbf{1}_{\calO_v^2})
  =L_v(s,\tilde\pi_v\times\pi_v)
\]
when $\vol(N_v A_v K_v)=\vol(K_v)=1$ \cite[Proposition~(2.3), pp.~511--512]{JS81}.
This is compatible with the Casselman--Shalika normalization used for F-1
\cite[Theorem~5.4, p.~227]{CS80}.

This input covers only the unramified places.  It does not specify the values,
signs, or archimedean normalization of the local correction factors required in
the global residue formula.

\subsection{Global integral and pole criterion}

Jacquet--Shalika construct the corresponding global Whittaker integral, prove
factorization into local integrals for pure tensors, and prove that it is
meromorphic and holomorphic for $\Re s>1$.  If the Fourier transform of the
global test function is nonzero at zero and $\pi'=\tilde\pi$, the integral has a
simple pole at $s=1$; conversely, such a pole forces $\pi'=\tilde\pi$
\cite[Section~4.5--4.7, especially Lemma~(4.6), pp.~550--552]{JS81}.

This establishes existence of a simple pole for suitable data, but not the
explicit positive residue formula desired in F-2.  A complete proof must:
\begin{enumerate}
  \item choose the global data and normalize all Haar measures;
  \item compute each factor in the finite bad-place product;
  \item prove that every relevant local factor is nonzero at $s=1$ and has the
        sign claimed in the final formula;
  \item identify the global constant with a positive norm-square;
  \item pass from $L(s,\tilde\pi\times\pi)$ to $\zeta_F(s)L(s,\pi,\Ad)$,
        including the bad local definitions and possible central-character
        twists.
\end{enumerate}

\subsection{Adjoint value}

The current proof attempt required the exact input
\[
  L(1,\pi,\Ad)>0.
\]
The ledger's source check of Shahidi's 1981 Theorems 5.2--5.3 found
nonvanishing of pair $L$-functions and of symmetric third/fourth powers, not
this adjoint statement.  The pair statement at $t=0$ concerns a product with
the zeta factor and therefore does not by itself separate the zeta pole from a
zero of the adjoint factor.

Accordingly, the following bridge remains open: either locate a primary theorem
that states $L(1,\pi,\Ad)>0$ under the required hypotheses, or prove a complete
bridge from the Jacquet--Shalika simple-pole criterion to nonvanishing of
$L(1,\pi,\Ad)$, then determine its sign from real-axis normalization and
continuity.  Until that is done, this input is unavailable.

\section{What must not be inferred}

Even after F-2 is proved, it would be pointwise in $\pi$.  It would not provide
a lower bound uniform in the level, a derivative/logarithmic-derivative bound,
or exclusion of a real zero approaching $s=1$.  Those quantitative implications
remain separate obligations in proof documents 03 and 04.

\begin{thebibliography}{9}

\bibitem{CS80} W.~Casselman and J.~Shalika, The unramified principal series of
  $p$-adic groups II: the Whittaker function, Compositio Math.\ \textbf{41}
  (1980), 207--231.

\bibitem{JS81} H.~Jacquet and J.~A.~Shalika, On Euler products and the
  classification of automorphic representations I, Amer.\ J.\ Math.\ \textbf{103}
  (1981), 499--558.

\end{thebibliography}

\end{document}
